\documentclass[11pt]{article}

% --- encoding & fonts (first) ---
\usepackage[T1]{fontenc}
\usepackage{lmodern}

% --- math stack ---
\usepackage{amsmath}
\usepackage{amssymb}
\usepackage{mathtools}      % supersets + fixes amsmath
\usepackage{amsthm}
\usepackage{bm}

% --- tables, graphics, colour, lists, quotes ---
\usepackage{booktabs}
\usepackage{graphicx}
\usepackage{xcolor}
\usepackage{enumitem}
\usepackage{csquotes}

% --- language & layout ---
\usepackage[american]{babel}
\usepackage[margin=1in]{geometry}

% --- typography (after the font) ---
\usepackage{microtype}
\setlength{\emergencystretch}{3em}

% --- references: hyperref near-last, xurl after, cleveref last ---
\usepackage[colorlinks=true,linkcolor=blue,citecolor=blue,urlcolor=blue,breaklinks=true]{hyperref}
\usepackage{xurl}
\usepackage[capitalize,noabbrev]{cleveref}   % only needed if you use \cref; harmless otherwise

% --- theorem environments: ONE section-based shared counter (retrievable refs: Lemma 2.1, Theorem 4.1) ---
\newtheorem{theorem}{Theorem}[section]
\newtheorem{proposition}[theorem]{Proposition}
\newtheorem{lemma}[theorem]{Lemma}
\newtheorem{corollary}[theorem]{Corollary}
\newtheorem{conjecture}[theorem]{Conjecture}
\newtheorem{question}[theorem]{Question}
\theoremstyle{definition}
\newtheorem{definition}[theorem]{Definition}
\theoremstyle{remark}
\newtheorem{remark}[theorem]{Remark}

% --- project macros ---
\newcommand{\R}{\mathbb{R}}
\newcommand{\sgn}{\operatorname{sign}}

\title{Averaged-Jacobian failure at the first open width of the Neural Jacobian Conjecture\thanks{Prompted by the Neural Jacobian Conjecture of Chen and Jiang, arXiv:2606.10806.}}
\author{Logan M. Dixon\thanks{Independent Researcher. Email:
\texttt{lmdixon23@gmail.com}. ORCID: 0009-0001-0592-462X.}}
\date{July 2026}

\begin{document}
\maketitle

\begin{abstract}
For sigmoid ridge networks $F(x)=A\sigma(Bx+c)$, the Neural Jacobian Conjecture asks whether
$\det DF>0$ on $\mathbb R^n$ forces global injectivity. At widths $N=n$ and $N=n+1$, the known
proofs reduce to mean-value univalence: the determinant is affine in the single varying slot and
therefore commutes with segment averaging. We show that this mechanism fails at the first open
width. An explicit network with $n=2$ and $N=4$ satisfies $\det DF(x)>0$ for every $x$, while one
segment-averaged Jacobian has negative determinant; continuity then gives an intermediate pair with
singular averaged Jacobian. Thus pointwise Jacobian positivity does not imply the nonsingularity
condition used by the mean-value certificate. Global positivity is established by a directional tail
bound and an exhaustive compact branch-and-bound certificate, whereas the averaged determinant is
evaluated from exact sigmoid divided differences. A sign-vector duality argument excludes fourteen
of the sixteen possible collision-direction patterns, confining any collision to two antipodal open
sectors. The witness's injectivity remains open, as does the conjecture for $n\ge2$ and
$N\ge n+2$. The example isolates the obstruction: once the kernel dimension reaches two, the
determinant becomes genuinely multilinear and need not survive averaging.
\end{abstract}

\medskip
\noindent\textbf{2020 Mathematics Subject Classification.} 26B10 (primary); 15A15, 15B35, 65G20, 68T07.

\smallskip
\noindent\textbf{Keywords.} global univalence; averaged Jacobian; Neural Jacobian Conjecture;
sigmoid ridge network; validated numerics; sign-vector condition.

\section{Introduction}\label{sec:intro}
For a sigmoid ridge network
\[
F(x)=A\sigma(Bx+c),\qquad
A\in\mathbb R^{n\times N},\quad B\in\mathbb R^{N\times n},\quad c\in\mathbb R^N,
\]
let $\mathcal N_{n,N}$ denote this class and let $\mathcal N^+_{n,N}$ be its subclass with
$\det DF(x)>0$ for every $x$. The Neural Jacobian Conjecture asks whether every
$F\in\mathcal N^+_{n,N}$ is one-to-one. The conjecture is known at $N=n$ and $N=n+1$. Its first unresolved width is
$N=n+2$, beginning with $n=2$, $N=4$.

The closed widths share a mean-value mechanism. A collision produces a segment-averaged Jacobian,
and when the kernel dimension is one the determinant is affine in the only varying slot. Averaging
therefore preserves its sign, so the averaged Jacobian remains nonsingular. At kernel dimension two,
the determinant is genuinely multilinear and this commutation need not survive.

The main result constructs that failure inside the conjecture's class. For an explicit
$2$-input, $4$-ridge network, the pointwise determinant is positive on all of $\mathbb R^2$, yet the
averaged determinant is negative for one segment and vanishes for an intermediate pair. The example
does not refute the Neural Jacobian Conjecture: it separates pointwise positivity from the classical
mean-value certificate that proves the smaller-width cases.

The construction has two mathematical layers. Cauchy--Binet expresses both pointwise and averaged
determinants through the same mixed-minor polynomial, evaluated respectively on derivative weights
and their segment averages. A unique negative minor allows an average to cross the determinant-zero
locus even though the pointwise derivative-weight path stays positive. A tail domination argument
and a finite compact certificate establish global pointwise positivity. Sign-vector duality then
rules out fourteen of the sixteen possible collision directions, leaving two antipodal sectors and a
concrete open injectivity problem.

Thus the obstruction is not local orientation but alignment: whenever an averaged Jacobian is
singular, the displacement vector must avoid its kernel. The final sections formulate this remaining
condition and separate the universal theorem-level claims from the platform-dependent frontier study.

\section{Divided differences and averaged Jacobians}\label{sec:setting}

Fix $F\in\mathcal N_{n,N}$ and write
$\sigma(t)=(1+e^{-t})^{-1}$. Let $b_i^{\mathsf T}$ denote the $i$-th row of
$B$, let $a_i$ denote the $i$-th column of $A$, and set
$t_i(x)=b_i\cdot x+c_i$. The subclass $\mathcal N^+_{n,N}$ consists of the
networks with $\det DF(x)>0$ for every $x$.

\begin{lemma}[Divided-difference identity]\label{lem:dd}
For $p,q\in\R^n$ with $u=q-p$,
\[
F(q)-F(p)=A\,\Delta(p,q)\,B\,u,
\qquad
\Delta(p,q)=\operatorname{diag}\bigl(\delta_1,\dots,\delta_N\bigr),
\]
where $\delta_i$ is the divided difference of $\sigma$ along the $i$-th
preactivation,
$\delta_i=\bigl(\sigma(t_i(q))-\sigma(t_i(p))\bigr)/\bigl(t_i(q)-t_i(p)\bigr)$
when $t_i(q)\ne t_i(p)$ and $\delta_i=\sigma'(t_i(p))$ otherwise. Moreover
$\Delta(p,q)=\int_0^1 D\bigl(p+tu\bigr)\,dt$, where
$D(x)=\operatorname{diag}\bigl(\sigma'(t_i(x))\bigr)$, so that
$A\,\Delta(p,q)\,B=\overline{DF}(p,q)$ is precisely the segment-averaged
Jacobian. In particular every $\delta_i$ lies in $(0,\tfrac14]$.
\end{lemma}

\begin{proof}
Componentwise, $\sigma(t_i(q))-\sigma(t_i(p))=\delta_i\,(b_i\cdot u)$ by
the definition of the divided difference, and
$\delta_i=\int_0^1\sigma'\bigl(t_i(p)+s\,b_i\cdot u\bigr)\,ds$ by the
fundamental theorem of calculus; the integral form also covers the case
$b_i\cdot u=0$. Strict positivity of $\delta_i$ follows from
$\sigma'>0$, and $\delta_i\le\sup\sigma'=\tfrac14$. Since
$DF(x)=A\,D(x)\,B$, linearity of the integral gives
$A\Delta B=\int_0^1 DF(p+tu)\,dt$.
\end{proof}

By the Cauchy--Binet formula, for any positive diagonal
$\Delta=\operatorname{diag}(\delta)$,
\begin{equation}\label{eq:cb}
\det\bigl(A\,\Delta\,B\bigr)
=\sum_{|I|=n} m_I\,\delta_I,
\qquad
m_I:=\det(A_{\cdot I})\,\det(B_{I\cdot}),
\quad
\delta_I:=\prod_{i\in I}\delta_i,
\end{equation}
the sum over all $n$-element subsets $I\subseteq\{1,\dots,N\}$. We call the
$m_I$ the \emph{minor products} of the pair $(A,B)$, and we call $(A,B)$
\emph{mixed-minor} if the nonzero $m_I$ do not all share one sign.

\section{The width-\texorpdfstring{$(n{+}1)$}{(n+1)} case is mean-value univalence}
\label{sec:comment}

The width-$(n+1)$ theorem of \cite{moonshine} is an immediate consequence of classical
mean-value univalence. The same argument identifies the precise mechanism, corrects a sign
normalization in the algebraic proofs, and fixes the scope of the open regime.

\begin{lemma}[Averaged-Jacobian criterion; folklore, cf.\
\cite{mcleod,parthasarathy,fmr}]\label{lem:mv}
Let $\Omega\subseteq\R^n$ be open and convex and $G\in C^1(\Omega,\R^n)$.
If for all $x,y\in\Omega$ the averaged Jacobian
$\overline{DG}(x,y):=\int_0^1 DG\bigl(y+t(x-y)\bigr)\,dt$ is nonsingular,
then $G$ is injective.
\end{lemma}
\begin{proof}
$G(x)-G(y)=\overline{DG}(x,y)\,(x-y)$ by the fundamental theorem of
calculus applied componentwise along the segment, which lies in $\Omega$ by
convexity. Nonsingularity forces $x=y$ whenever $G(x)=G(y)$.
\end{proof}

\begin{proposition}[Subsumption of {\cite[Lemma~3.3]{moonshine}}]
\label{prop:subsume}
Let $\Omega\subseteq\R^n$ be open and convex, $h\in C^1(\Omega)$,
$L:\R^{n+1}\to\R^n$ linear, and $T(y)=L(y,h(y))$ with $\det DT$
nonvanishing on $\Omega$. Then $T$ is injective.
\end{proposition}
\begin{proof}
Since $DT(y)=L\,(\mathrm{Id},Dh(y))$ and $\det DT$ is nonvanishing,
$L$ has rank $n$; hence $\ker L$ is one-dimensional. As in
\cite{moonshine}, after an invertible
target change and, in the second case, an invertible source change, $T$
takes one of the two normal forms
$T(y)=y+p\,h(y)$ or $\widetilde T(u,s)=(u,\widetilde h(u,s))$.
In the first form $\det DT(y)=1+\nabla h(y)\cdot p$ is \emph{affine} in
$\nabla h$, so
$\det\overline{DT}(x,y)=1+\bigl(\int_0^1\nabla h\bigr)\cdot p
=\int_0^1\det DT\,dt$.
In the second form $\det D\widetilde T=\partial_s\widetilde h$ is affine in
the varying slot, and again
$\det\overline{D\widetilde T}=\int_0^1\det D\widetilde T\,dt$.
In both cases $\det DT$ is continuous, nonvanishing on the connected set
$\Omega$, hence of constant sign, so the integral is nonzero. Lemma~\ref{lem:mv}
applies.
\end{proof}

\begin{corollary}[{\cite[Theorem~3.2]{moonshine}}]
Every $F\in\mathcal N^{+}_{n,n+1}$ is injective.
\end{corollary}

\begin{remark}[Mechanism]\label{rem:mechanism}
The proof of Proposition~\ref{prop:subsume} works because, at kernel
dimension one, the determinant is affine in the single varying slot, so
\emph{determinant commutes with segment averaging}. At kernel dimension
two or more the determinant is genuinely multilinear and the commutation
can fail; Section~4 exhibits a certified failure inside the sigmoid ridge
class itself. This, rather than fiber dimension per se, is the boundary
between the closed and open regimes of the Neural Jacobian Conjecture.
\end{remark}

\begin{remark}[Sign repair for the algebraic proofs of \cite{moonshine}]
Both algebraic proofs normalize the first $n$ preactivations by an affine
change $y=B_1x+c_1$. Since
$\det D(F\circ B_1^{-1})(y)=\det DF(x)\,\det B_1^{-1}$,
positivity need not survive when $\det B_1<0$; the arguments use only
constant nonzero sign, which does survive. The geometric proof of
\cite{moonshine} is unaffected.
\end{remark}

\begin{remark}[Scope of the open case]
For $n=1$ and every width $N$, any $F\in\mathcal N^{+}_{1,N}$ satisfies
$F'>0$ and is strictly increasing, hence injective. The unresolved regime
of the conjecture is therefore $n\ge 2$, $N\ge n+2$, not $N\ge n+2$
unqualified.
\end{remark}

\section{A hierarchy of injectivity certificates}\label{sec:hierarchy}

Identity \eqref{eq:cb} organizes three certificates of increasing
locality. We state them for the sigmoid ridge class; the first two are
specializations of general frameworks
\cite{mhr,muelleretal,fmr}.

\begin{proposition}[Fixed-sign minors certify the whole family]
\label{prop:c1}
If the minor products $m_I$ of $(A,B)$ are not all zero and the nonzero
ones all share one sign, then
\emph{every} map of the form $x\mapsto Af(Bx+c)$ built on $(A,B)$, where
$f$ is applied componentwise and each component is strictly increasing
and differentiable, is injective --- for every $c$ and every such $f$.
\end{proposition}
\begin{proof}
By Lemma~\ref{lem:dd} a collision $F(p)=F(q)$, $u=q-p\ne0$, forces
$A\Delta Bu=0$ with $\Delta$ a positive diagonal. By \eqref{eq:cb} every
positive diagonal makes $\det(A\Delta B)$ a sum of terms of one sign, not
all zero (some $m_I\ne0$ and all $\delta_I>0$), hence nonzero; so
$A\Delta B$ is nonsingular and $u=0$. The argument never uses the specific
activation: the identity of Lemma~\ref{lem:dd} is definitional
slot-by-slot, divided differences of a strictly increasing component are
strictly positive, and slots with $b_i\cdot u=0$ contribute nothing to
$\Delta Bu$, so their diagonal entries may be chosen positive. This is the
sigmoid-ridge instance of the sign-vector injectivity conditions of
\cite{mhr,muelleretal}; in the terminology of \cite{mhr}, it is the case
$\sgn(\ker A)\cap\sgn(\operatorname{im}B)=\{0\}$ certified through minors
(cf.\ \cite[Cor.~4]{mhr}).
\end{proof}

\begin{proposition}[Segment-average positivity certifies one map]
\label{prop:c2}
If $\det\overline{DF}(p,q)\ne0$ for all $p,q$, then $F$ is injective.
\end{proposition}
\begin{proof}
Lemma~\ref{lem:mv} with $G=F$, $\Omega=\R^n$; equivalently,
Lemma~\ref{lem:dd} plus nonsingularity of $A\Delta(p,q)B$.
\end{proof}

The third condition, membership in $\mathcal N^{+}$ (pointwise
positivity), is exactly the hypothesis of the NJC. Whether it certifies
injectivity for $n\ge2$, $N\ge n+2$ is the open question. The established
implications are
\[
\text{fixed-sign minors}
\;\Longrightarrow\;
\text{averaged-Jacobian nonsingularity}
\;\Longrightarrow\;
\text{injectivity}.
\]
By contrast, the unresolved implication is
\[
\mathcal N^{+}\;\overset{?}{\Longrightarrow}\;\text{injectivity}.
\]
Section~\ref{sec:witness} shows that the conjectural arrow cannot, in
general, be proved by inserting averaged-Jacobian nonsingularity as an
intermediate certificate: a mixed-minor $\mathcal N^{+}$ network can have
a singular averaged Jacobian. For fixed-sign networks the insertion does
succeed, by Proposition~\ref{prop:c1}.

\begin{remark}[A convex-gradient subfamily]\label{rem:gradient}
When $B$ has rank $n$, families with $A=RB^{\mathsf T}$, $R$ symmetric
positive definite, are
gradients of strictly convex potentials composed with a fixed linear
isomorphism, hence injective for reasons independent of the conjecture;
their minor products are automatically of one sign. Probing the open
regime requires mixed-minor pairs, which is where our search
(Section~\ref{sec:frontier}) operates.
\end{remark}

\section{Separation at \texorpdfstring{$n=2$, $N=4$}{n=2, N=4}}
\label{sec:witness}

\subsection{The witness}

Let $F(x)=A\sigma(Bx+c)$ with
\[
A=\begin{pmatrix}
-0.2398938710 & \phantom{-}0.2126956751 & -1.8416162379 & -0.0023672752\\
-0.7053370168 & \phantom{-}1.2188352888 & -2.8603622059 & \phantom{-}0.0955954101
\end{pmatrix},
\]
\[
B=\begin{pmatrix}
-0.2535883378 & -0.1878587713\\
\phantom{-}0.1326989170 & \phantom{-}1.0232004273\\
-1.9014473780 & \phantom{-}1.0928150941\\
-0.3810233277 & -0.1438462106
\end{pmatrix},
\qquad
c=\begin{pmatrix}
-0.9291312391\\ \phantom{-}0.1478144974\\ -0.9119451690\\ \phantom{-}0.0633145878
\end{pmatrix},
\]
displayed here to ten decimals; the repository file
\texttt{data/witness.json} carries the canonical full-precision values,
and every constant below is machine-computed from that file. The six minor
products are
\[
m_{12}=0.0333916540,\quad
m_{13}=0.3887018928,\quad
m_{14}=0.0008635680,
\]
\[
m_{23}=3.4206863493,\quad
m_{24}=0.0086086719,\quad
m_{34}=-0.1261291707,
\]
so $(A,B)$ is mixed-minor with a unique negative product.

\subsection{Global positivity}

\begin{theorem}[Global positivity]\label{thm:cert}
$\det DF(x)>0$ for all $x\in\R^2$; that is, $F\in\mathcal N^{+}_{2,4}$.
\end{theorem}

The proof has a tail part and a compact part. Outside the disk of radius
$R_0=601$, the unique negative term in \eqref{eq:cb} is dominated by the
positive terms; the logarithmic domination margin is at most $-0.008291$.
Inside the disk, an outward-rounded branch-and-bound covers the domain by
$92{,}696$ terminal boxes and gives a strictly positive determinant lower
bound on each box. The public certificate stores the subdivision and the
bounds, while a dependency-free checker recomputes every load-bearing
inequality and rejects malformed or ambiguous inputs.

\subsection{Failure of the averaged-Jacobian certificate}

\begin{theorem}[Separation]\label{thm:sep}
Let $p=(-2.76601947,\,12)$ and $q=(7.08288117,\,12)$. Then
\[
\det DF\bigl(p+t(q-p)\bigr)\;\ge\;10^{-6}
\quad\text{for all }t\in[0,1],
\qquad\text{while}\qquad
\det\overline{DF}(p,q)\;<\;0,
\]
the latter approximately $-1.7200\times10^{-5}$.
Moreover, writing $u=q-p$, the function
$s\mapsto\det\overline{DF}(p,\,p+su)$ is continuous on $(0,1]$ with
positive limit $\det DF(p)$ (approximately $1.4051\times10^{-5}$) as
$s\to0^{+}$, so there is
$s^{*}\in(0,1)$ --- numerically $s^{*}\approx0.7766094$ --- at which the
averaged Jacobian $\overline{DF}(p,\,p+s^{*}u)$ is itself
\emph{singular}. Consequently the hypothesis of Lemma~\ref{lem:mv}
(equivalently, of the mean-value univalence criteria of
\cite{mcleod,parthasarathy,fmr}) fails for $F$, although
$F\in\mathcal N^{+}_{2,4}$ with certified positivity; in addition the
segment hull
$\mathcal H:=\operatorname{conv}\{DF(p+t(q-p)):t\in[0,1]\}$ contains a
singular matrix.
\end{theorem}

\begin{proof}
An outward-rounded one-dimensional subdivision proves the displayed
bound $10^{-6}$ on the whole segment. For reference, the smallest
certified lower bound among $201$ sampled subintervals is
$1.2620\times10^{-6}$. The averaged determinant and the value at $p$ are
direct evaluations: Lemma~\ref{lem:dd} gives
\[
\overline{DF}(p,q)=A\Delta(p,q)B.
\]

The average of a continuous matrix path lies in the convex hull of its
values, so $\overline{DF}(p,q)\in\mathcal H$; this is the matrix
mean-value theorem of McLeod \cite{mcleod}, and also follows from the
finite convex-combination form in \cite[Prop.~4.3]{fmr}. Every matrix
$DF(p+t(q-p))$ has positive determinant, whereas
$\overline{DF}(p,q)$ has negative determinant. Continuity of $\det$ on a
line segment in $\mathcal H$ therefore gives a singular matrix in the
hull.

Finally, $\overline{DF}(p,p+su)$ depends continuously on $s$ by the
integral formula in Lemma~\ref{lem:dd}. Its determinant tends to
$\det DF(p)>0$ as $s\to0^+$ and is negative at $s=1$. The intermediate
value theorem gives $s^*\in(0,1)$; bisection yields
$s^*=0.7766094$ to seven digits, with second point
$\approx(4.8827295,12)$.
\end{proof}

\begin{remark}[What is, and is not, new here]\label{rem:fmr45}
Failure of the averaged-Jacobian certificate for a map with positive
pointwise determinant is classical \emph{when the map is not injective}:
by the contrapositive of Lemma~\ref{lem:mv}, every non-injective example
exhibits it, a standard example in the classical global univalence
literature \cite{galenikaido,parthasarathy} being the complex exponential
$G(x,y)=e^{x}(\cos y,\sin y)$ with $\det DG=e^{2x}>0$ and period
$2\pi$ in $y$. The content of Theorem~\ref{thm:sep} is therefore not hull or average
degeneracy by itself, but its occurrence under global pointwise positivity
at the first open width. The witness is not known to be non-injective; if
the conjecture holds, it separates the mean-value certificate from the
truth of injectivity. Within \cite{fmr}, the closest point of
contact is Example~4.5, which separates certificate classes for a
monomial map ($Q(D_G)$ contains singular matrices while a specialized
class $q(B)$ certifies injectivity) and leaves the singularity status of
$\operatorname{conv}(J_G)$ unaddressed; Theorem~\ref{thm:sep} exhibits a
singular averaged Jacobian itself. To our knowledge this is the first
certified instance in the sigmoid ridge class.
\end{remark}

\subsection{Partial injectivity of the witness: sector exclusion}

Theorem~\ref{thm:sep} removes the standard route to injectivity for the
witness, and raises the natural question of whether the witness is
injective at all. We do not know; what we can prove confines any
hypothetical collision sharply. Adopting the sign-vector viewpoint of
\cite{mhr}: by Lemma~\ref{lem:dd}, a collision with direction $u$ forces
$\gamma:=\Delta Bu\in\ker A\setminus\{0\}$ with
$\sgn(\gamma)=\sgn(Bu)$ componentwise.

\begin{theorem}[Sector exclusion]\label{thm:sector}
The sign vector $\sgn(Bu)$, $u\ne0$, takes exactly $16$ values: $8$
open-sector patterns and $8$ boundary-ray patterns. For $14$ of them
there is a certificate $\lambda\in\R^2$ with
$\sigma_i\,(A^{\mathsf T}\lambda)_i>0$ for every $i$ with $\sigma_i\ne0$,
and then no nonzero $\gamma\in\ker A$ realizes the pattern: indeed
$0=\lambda^{\mathsf T}A\gamma=\sum_i(A^{\mathsf T}\lambda)_i\gamma_i$
would be a sum of nonnegative terms, at least one strictly positive. The
two remaining patterns $\pm(-,+,+,-)$ are realized in $\ker A$, so the
exclusion is sharp. Consequently any collision direction of the witness
lies strictly inside the open sector
$S^{+}=(60.1128^\circ,\,110.6828^\circ)$ or its antipode; in particular
$F$ is injective along every line whose direction lies outside
$S^{+}\cup S^{-}$, including all eight sector boundaries.
\end{theorem}

The fourteen dual certificates have minimum margin $0.017442$ after
normalizing $\|\lambda\|=1$; each is checked by one matrix-vector
product. The witness segment direction $(1,0)$ has pattern $(-,+,-,-)$
and is excluded, consistent with Theorem~\ref{thm:sep}.

Beyond direction, magnitude is also confined: writing $\delta_{(i)}$
for the divided difference of unit $i$, a collision makes the averaged
Jacobian singular along the collision direction, so the Cauchy--Binet
sum \eqref{eq:cb} vanishes there; moving the unique negative term
across, discarding the remaining nonnegative terms, dividing by
$\delta_{(3)}>0$, and bounding $\delta_{(4)}\le\tfrac14$ yields the
linear domination inequality
$m_{23}\,\delta_{(2)}+m_{13}\,\delta_{(1)}\le|m_{34}|/4$, numerically
$3.4206863\,\delta_{(2)}+0.3887019\,\delta_{(1)}\le 0.0315323$, hence
$\delta_{(2)}\le 0.0092182$ and $\delta_{(1)}\le 0.0811221$, and thus
saturated preactivations $|t_2|\ge3.300$ and $|t_1|\ge1.1255$ somewhere
on the collision segment. All unit indices in this paragraph are 1-based.

\begin{question}[Injectivity of the witness]\label{q:witness}
Is the witness network injective? Equivalently, does the finite-dimensional confined system
--- collision direction strictly inside $S^{+}\cup S^{-}$, the magnitude
and saturation bounds above, and $\Delta Bu\in\ker A$ --- admit a solution?
\end{question}

\section{Frontier study}\label{sec:frontier}

We also tested how often the phenomenon in Theorem~\ref{thm:sep}
appeared in the search distribution used to find the witness. We
sampled random mixed-minor candidates at $n=2$, $N=4$ using the pinned
seeds in \texttt{code/njc\_v3.py}. In the study environment, $911$
draws produced $45$ accepted mixed-minor networks certified in
$\mathcal N^{+}$; for $44$ of them segment-average positivity was also
certified over the searched pair space, while exactly one, the witness,
exhibited a negative segment average.

The pinned seeds fix the random draws, but not every floating-point
accept/reject decision made by the search code. These counts should
therefore be read as the recorded output of the study environment
rather than as platform-independent enumerative data. The accepted
networks ship in \texttt{data/e3.pkl}. The witness itself, and every
load-bearing certified claim about it, are platform-independent: they
are anchored to \texttt{data/witness.json} and verified by the
dependency-free checker.

No collision was found in the searched region for any accepted
network. A second implementation under an independent sampling protocol
also accepted $45$ networks and found no additional frontier instance.
Across the two runs, one frontier instance appeared among roughly
ninety accepted outputs; this is a search-distribution observation, not
an estimate of prevalence in the class.

Two caveats are important. First, the acceptance and certification
filters impose magnitude thresholds, so borderline cases may be
underrepresented. Second, the absence of detected collisions is
evidence only about the searched regions, not a certificate of
injectivity.

\section{The Neural Jacobian Conjecture as kernel alignment}
\label{sec:njc}

By Lemma~\ref{lem:dd}, injectivity of $F$ is exactly the statement
\[
u\ne0\ \Longrightarrow\ A\,\Delta(p,p+u)\,B\,u\ne0 .
\]
The closed cases prove this by making $A\Delta B$ nonsingular for every
positive diagonal arising as a divided-difference matrix.
Theorem~\ref{thm:sep} shows that this route already fails as a general
proof strategy at the first open width: for the certified $n=2$, $N=4$
witness, $\det\bigl(A\Delta(p,q)B\bigr)$ changes sign as $(p,q)$
varies, so there exist pairs at which the averaged Jacobian is singular.
What remains of the conjecture is therefore a genuine alignment
statement: whenever $A\Delta(p,q)B$ is singular, the direction $q-p$ must
avoid its kernel. Our witness makes this concrete --- its kernel-avoidance
is exactly Question~\ref{q:witness} --- and the sector-exclusion method
suggests that sign-vector techniques \cite{mhr,muelleretal,fmr}, rather
than averaging, are the natural tools for the open regime. The conjecture
itself we neither prove nor refute; the frontier study is consistent
with it, and the separation example
explains why it has resisted mean-value methods.

Thus the first open width already separates pointwise positivity from
the averaged-Jacobian certificate. The Neural Jacobian Conjecture, if
true, must rely on a finer alignment constraint between
divided-difference kernels and displacement directions, rather than on
nonsingularity of the averaged Jacobian alone.

\section{Verification and reproducibility}\label{sec:repro}

\paragraph{Repository and command.}
The public repository is
\url{https://github.com/lmdixon23/njc-separation}. Install the locked dependencies with
\texttt{python -m pip install -r requirements-lock.txt}, then run
\texttt{bash run\_all.sh}. The suite verifies the SHA-256 manifest before running the mathematical
and certificate checks.

\paragraph{Proof and certificate boundary.}
The divided-difference identities, width-$(n+1)$ argument, certificate hierarchy, and sign-vector
implications are proved analytically. Global positivity of the explicit witness uses the stored
compact-and-tail certificate checked by \texttt{CHECK.py}. The segment lower bound uses
outward-rounded interval subdivision. The frontier counts describe the shipped search run and are
not asserted as platform-independent enumerative data.

\paragraph{Independent reconstruction.}
\texttt{BLIND\_CHECK.py} reads only the canonical witness data and independently recomputes the mixed
minors, negative averaged determinant, and positive pointwise samples on the displayed segment. The
proof-bearing checker does not trust verdict fields in the certificate; it reconstructs the required
bounds and partition conditions.

\paragraph{Limitations.}
The witness's injectivity is open. The sector theorem supplies necessary conditions for a collision,
not a collision exclusion. The Neural Jacobian Conjecture remains open for $n\ge2$ and
$N\ge n+2$, and the frontier study does not estimate prevalence in that class.

\section*{Declarations}
\paragraph{Funding.} No funding was received for this work.
\paragraph{Competing interests.} The author declares no competing interests.
\paragraph{Data and code availability.} The canonical witness, certificate, search data, SHA-256
manifest, and verification scripts are available at
\url{https://github.com/lmdixon23/njc-separation}.
\paragraph{AI assistance.} Text-to-text generative AI tools, including Claude and GPT, assisted with
candidate exploration, code drafting and review, adversarial checking, and editorial revision. Any
AI-assisted mathematical claim, source attribution, code fragment, constant, or numerical output
retained in the work was checked independently of the AI response through hand derivation, source
verification, exact recomputation, or validated numerical computation as appropriate. The author
determined the research direction and acceptance criteria and assumes full responsibility for the
work.

\begin{thebibliography}{9}
\bibitem{moonshine} X.~Chen and X.~Jiang, \emph{Moonshine: an autonomous
mathematical research agent centered on conjecture generation},
arXiv:2606.10806, 2026.
\bibitem{galenikaido} D.~Gale and H.~Nikaido, \emph{The Jacobian matrix and
global univalence of mappings}, Math.\ Ann.\ 159 (1965), 81--93.
\bibitem{mcleod} R.~M.~McLeod, \emph{Mean value theorems for vector valued
functions}, Proc.\ Edinburgh Math.\ Soc.\ 14 (1965), 197--209.
\bibitem{parthasarathy} T.~Parthasarathy, \emph{On Global Univalence
Theorems}, Lecture Notes in Mathematics 977, Springer, 1983.
\bibitem{fmr} E.~Feliu, S.~M\"uller, and G.~Regensburger,
\emph{Characterizing injectivity of classes of maps via classes of
matrices}, Linear Algebra Appl. 580 (2019), 236--261. arXiv:1901.11272.
\bibitem{muelleretal} S.~M\"uller, E.~Feliu, G.~Regensburger, C.~Conradi,
A.~Shiu, and A.~Dickenstein, \emph{Sign conditions for injectivity of
generalized polynomial maps with applications to chemical reaction
networks and real algebraic geometry}, Found.\ Comput.\ Math.\ 16 (2016),
69--97.
\bibitem{mhr} S.~M\"uller, J.~Hofbauer, and G.~Regensburger, \emph{On the
bijectivity of families of exponential/generalized polynomial maps},
SIAM J.\ Appl.\ Algebra Geom.\ 3 (2019), no.~3, 412--438.
arXiv:1804.01851.
\end{thebibliography}
\end{document}
